\documentclass[journal]{IEEEtran}

% *** PACKAGES ***
\usepackage{amsmath,amssymb,amsfonts}
\usepackage{graphicx}
\usepackage{cite}
\usepackage{url}
\usepackage{multirow}
\usepackage{booktabs}
\usepackage{amsthm}

% *** THEOREM ENVIRONMENTS ***
\newtheorem{theorem}{Theorem}[section]
\newtheorem{proposition}[theorem]{Proposition}
\newtheorem{lemma}[theorem]{Lemma}
\newtheorem{corollary}[theorem]{Corollary}
\theoremstyle{definition}
\newtheorem{definition}[theorem]{Definition}
\newtheorem{remark}[theorem]{Remark}

\newcommand{\coloneqq}{\triangleq}

\begin{document}

\title{A Spectral Framework for Measuring Output-Space Irreversibility in Machine Unlearning}

\author{%
Yan Zhang\\
\IEEEauthorblockA{%
College of Mathematical Sciences,\\
Zhejiang University of Technology,\\
Hangzhou, China}%
\thanks{This work was done while the author was with the College of Mathematical Sciences, Zhejiang University of Technology.}%
\thanks{Manuscript received Month xx, 2026. This work was supported by Project No.~2026FRG05004 under the 2026 Central Government--Guided Local Science and Technology Development Fund. The implementation is available at \protect\url{https://github.com/EurusDemerzel/rii_unlearning}.}
\and
Chenggang Lu\\
\IEEEauthorblockA{%
College of Mathematical Sciences,\\
Zhejiang University of Technology,\\
Hangzhou, China}%
\thanks{Corresponding author: Chenggang Lu (e-mail: luchenggang@zjut.edu.cn).}
}

\markboth{IEEE Transactions on Information Forensics and Security}%
{Author \MakeLowercase{\textit{et al.}}: A Spectral Framework for Output-Space Irreversibility in Machine Unlearning}

\maketitle

\begin{abstract}
We present a spectral framework for quantifying output-space irreversibility in machine unlearning. The central object is a $2\times C$ \emph{empirical confusion matrix} whose rows are the averaged softmax predictions on the forget set $D_f$ and retain set $D_r$. We define the \textbf{Residual Information Index (RII)} $\rho = 1 - \sigma_1^2/(\sigma_1^2+\sigma_2^2)$ from its singular values, measuring the deviation from the rank-one ideal. We prove that $\rho=0$ is necessary and sufficient for \emph{output distributional irreversibility}: the forget-set membership indicator $X$ and the observed model output $Y$ are conditionally independent given $\theta'$, i.e., $I(D_f; Y \mid \theta') = 0$. This guarantee is restricted to the output layer---it does not claim parameter-level white-box security---but is precisely what matters for black-box GDPR compliance auditing. For class-level forgetting, where the standard RII suffers from known/unknown class asymmetry, we introduce the \textbf{Multi-Held-Out Projection Residual (MHPR)} $\rho_H$, which projects the forgotten class mean onto the subspace spanned by multiple held-out unseen classes, and we establish its theoretical guarantees (monotonic improvement, finite-sample confidence, and bias--variance decomposition). We further propose temperature-scaled MHPR to overcome subspace collapse in low-accuracy regimes. We provide (i) a finite-sample confidence bound with explicit subsampling guarantees, (ii) an $O(\eta^2)$ leakage bound for gradient-based unlearning, (iii) a spectral mutual information bound $I(X;Y) \le O(\rho)$, and (iv) a unified leakage decomposition $I(X;R) \le C_1\alpha^2 + C_2\eta^2$ separating direct signal leakage from spectral rank deficiency. Experiments on MNIST, Fashion-MNIST, CIFAR-10, CIFAR-100, and Tiny ImageNet (200 classes) across multiple model architectures confirm $\rho < 2 \times 10^{-3}$ under standard training, surging by $10\times$ under deliberate overfitting and $180\times$ under targeted class forgetting. MHPR with $K\ge 3$ held-out classes achieves $\rho_H \approx 0.046$ on MNIST, and with temperature scaling ($T=50$) achieves $\rho_H=0.021$ on Fashion-MNIST (88.7\% accuracy), outperforming the single-class LOCO baseline by an order of magnitude. We reconcile the paradox of near-zero RII with 77--94\% MIA accuracy by showing they measure orthogonal leakage axes: MIA detects distributional shift ($I_{\mathrm{dir}}$), while RII certifies sample-level output indistinguishability ($I_{\mathrm{cov}}$).
\end{abstract}

\begin{IEEEkeywords}
Rank-one channels, information-theoretic security, machine unlearning, singular value decomposition, residual information index, multi-held-out projection, membership inference attack.
\end{IEEEkeywords}

%=============================================================================
\section{Introduction}
\label{sec:intro}
%=============================================================================

The right to erasure, codified in Article~17 of the GDPR, increasingly applies to machine learning models: upon request, a model must eliminate the influence of specific training data. This has spurred the field of \emph{machine unlearning}~\cite{bourtoule2021machine,nguyen2025survey}, which designs algorithms $\mathcal{U}$ that transform trained parameters $\theta_0$ into $\theta' = \mathcal{U}(\theta_0, D_f, D_r)$ with reduced dependence on the forget set $D_f$.

A central challenge is \emph{how to verify that unlearning has actually succeeded}. Existing evaluation methods either rely on membership inference attacks (MIA) that conflate train/test distributional shift with forget/retain leakage, or on differential privacy bounds that require invasive training procedures. What is missing is a lightweight, black-box certificate that directly quantifies whether a specific forget set remains distinguishable in the model's outputs.

Our framework fills this gap by observing that \textbf{output-space irreversibility reduces to a spectral condition}: the distribution of outputs on the forget set and the retain set must be conditionally independent of forget-set membership. When this holds, the corresponding $2\times C$ input-output channel has rank one, and $I(D_f; Y \mid \theta') = 0$.

\textbf{Scope.} We emphasize that our framework measures irreversibility strictly in the \emph{output space}---the observable predictions of the model. We do not claim parameter-level white-box security. This scope is precisely aligned with black-box GDPR compliance auditing, where an external regulator can query the model API but cannot inspect internal weights. For white-box guarantees, RII should be combined with internal representation audits (Section~\ref{sec:threat}).

Our framework abstracts this insight via the \textbf{$2\times C$ empirical confusion matrix} $\mathbf{M}(\theta')$, whose two rows are the averaged softmax predictions on $D_f$ and $D_r$. The singular values $\sigma_1 \ge \sigma_2 \ge 0$ of $\mathbf{M}$ encode the degree of residual dependence:

\begin{itemize}
\item $\sigma_2 = 0 \iff \mathrm{rank}(\mathbf{M}) = 1 \iff$ perfect removal;
\item $\sigma_2 > 0 \iff$ information leaks through the output distribution.
\end{itemize}

We define the \textbf{Residual Information Index (RII)} $\rho = 1 - \sigma_1^2/(\sigma_1^2+\sigma_2^2) \in [0, 0.5]$ as a normalized, scale-invariant measure of this deviation. For class-level forgetting, where the rows naturally differ due to known/unknown class asymmetry, we introduce the \textbf{Multi-Held-Out Projection Residual (MHPR)}, which measures the residual of the forgotten class mean after projection onto the subspace of multiple held-out unseen classes.

\textbf{Contributions:}
\begin{enumerate}
\item \emph{Output-Space Spectral Framework} (Sec.~\ref{sec:channel}): The $2\times C$ empirical channel, with rank-one condition $\sigma_2=0$ as the certificate of output distributional irreversibility for machine unlearning.
\item \emph{Output Indistinguishability Criterion} (Theorem~\ref{thm:perfect}): $\rho=0 \iff I(X;Y\mid\theta')=0$, i.e., $\rho=0$ is necessary and sufficient for black-box indistinguishability of forget and retain samples.
\item \emph{MHPR for Class-Level Forgetting} (Sec.~\ref{sec:mhpr}): A multi-held-out projection that overcomes the single-class bias of LOCO, providing an asymptotically correct reference for class-level unlearning.
\item \emph{Additional Theoretical Results} (Sec.~\ref{sec:theory}, Appendix~\ref{sec:appendix}): Finite-sample confidence with subsampling guarantees (Proposition~\ref{prop:confidence}), gradient leakage $O(\eta^2)$ (Theorem~\ref{thm:finetune}), spectral MI bound $I \le O(\rho)$ (Theorem~\ref{thm:mi_bound}), and unified leakage decomposition $I(X;R) \le C_1\alpha^2 + C_2\eta^2$ (Theorem~\ref{thm:unified}).
\item \emph{Comprehensive Empirical Validation} (Sec.~\ref{sec:experiments}): 50+ experimental configurations spanning MNIST, Fashion-MNIST, CIFAR-10, CIFAR-100, and Tiny ImageNet (200 classes) across four architectures, four unlearning methods, random and targeted class forgetting, gradient intensity sweeps, fine-tuning rebound, overfitting sensitivity, MHPR evaluation, and multi-seed statistical validation.
\item \emph{Orthogonal Leakage Decomposition} (Sec.~\ref{sec:mia}): Formal proof that MIA and RII measure independent axes---distributional shift vs.\ sample-specific output divergence---resolving their apparent paradox.
\item \emph{Beyond Output: Completing the Picture} (Sec.~\ref{sec:threat}): Decomposition $I(D_f;\theta') = I(D_f;Y\mid\theta') + I(D_f;\theta'\mid Y)$, showing how RII combined with internal representation audits (MMD) provides a path toward complete unlearning certification.
\end{enumerate}

%=============================================================================
\section{Problem Setup and Spectral Framework}
\label{sec:channel}
%=============================================================================

\subsection{Problem Setup}

Let $D = \{(x_i, y_i)\}_{i=1}^N$ be a training dataset, $\theta_0 = \mathcal{A}(D)$ the parameters obtained by training algorithm $\mathcal{A}$, $D_f \subset D$ the forget set, $D_r = D \setminus D_f$ the retain set, with $N_f = |D_f|$, $N_r = |D_r|$. An unlearning algorithm $\mathcal{U}$ produces $\theta' = \mathcal{U}(\theta_0, D_f, D_r)$. The goal is to quantify the residual dependence $I(D_f; \theta' \mid \theta_0)$.

Rather than operating in parameter space (where $I(X;\theta)$ is intractable for deep networks), we project onto the output space via the model's predictive distribution. This is both computationally tractable and information-theoretically principled, as any information about $D_f$ that influences decisions must manifest in the output distribution.

\paragraph{Two unlearning scenarios.}
We distinguish two fundamentally different forgetting regimes throughout this paper.

\noindent\textbf{Sample-level unlearning (default setting).} The forget set $D_f$ comprises individual training examples drawn uniformly from all classes. Because the retained set $D_r$ still contains many examples of every class, the model's strong generalization capacity may cause it to behave similarly on $D_f$ and $D_r$ even \emph{without} any explicit unlearning. This scenario is relevant for privacy requests targeting specific records (e.g., one user's images) but is the easier case for unlearning certification.

\noindent\textbf{Class-level unlearning.} The forget set $D_f$ comprises all examples of one or more entire semantic classes (e.g., all 6,000 images of the digit ``5'' in MNIST, or all 5,000 images of ``cat'' in CIFAR-10). After removal, the model has never seen this class during training; its predictions on class-$f$ samples must therefore come exclusively from the inductive biases of the architecture, not from memorized patterns. This is the more stringent test of unlearning, and it exposes structure that sample-level unlearning can obscure. As we shall see, class-level forgetting yields an RII $\rho \approx 0.18$---two orders of magnitude above the sample-level baseline---demonstrating the large dynamic range of the spectral metric.

\subsection{The Empirical Confusion Matrix}

Let $f_{\theta'}: \mathcal{X} \to \mathbb{R}^C$ be the logit function and $\mathbf{p}(x;\theta') = \mathrm{Softmax}(f_{\theta'}(x)) \in \Delta^{C-1}$ the probability vector over $C$ classes.

\begin{definition}[Empirical Confusion Matrix]
\label{def:M}
\begin{equation}
\mathbf{M}(\theta') =
\begin{bmatrix}
\boldsymbol{\mu}_f^\top \\[2pt]
\boldsymbol{\mu}_r^\top
\end{bmatrix}
\in \mathbb{R}^{2\times C},
\qquad
\begin{aligned}
\boldsymbol{\mu}_f &= \frac{1}{N_f}\sum_{x\in D_f} \mathbf{p}(x;\theta'),\\
\boldsymbol{\mu}_r &= \frac{1}{N_r}\sum_{x\in D_r} \mathbf{p}(x;\theta').
\end{aligned}
\end{equation}
\end{definition}
The matrix $\mathbf{M}$ is a $2\times C$ empirical channel: its two rows are the average predictive distributions on the forget and retain sets. If unlearning succeeds, these distributions should be statistically indistinguishable, implying $\boldsymbol{\mu}_f \approx \boldsymbol{\mu}_r$ and $\mathrm{rank}(\mathbf{M}) \approx 1$.

The SVD of $\mathbf{M}$ is particularly simple since it has only two rows:
\begin{equation}
\mathbf{M} = \sigma_1 \mathbf{u}_1 \mathbf{v}_1^\top + \sigma_2 \mathbf{u}_2 \mathbf{v}_2^\top,
\label{eq:svd_M}
\end{equation}
where $\sigma_1 \ge \sigma_2 \ge 0$, $\mathbf{u}_i \in \mathbb{R}^2$, $\mathbf{v}_i \in \mathbb{R}^C$. The squared singular values are the eigenvalues of the $2\times 2$ Gram matrix:
\begin{equation}
\mathbf{M}\mathbf{M}^\top =
\begin{bmatrix}
\|\boldsymbol{\mu}_f\|^2 & \boldsymbol{\mu}_f^\top\boldsymbol{\mu}_r \\[2pt]
\boldsymbol{\mu}_f^\top\boldsymbol{\mu}_r & \|\boldsymbol{\mu}_r\|^2
\end{bmatrix}.
\label{eq:gram}
\end{equation}

When $\|\boldsymbol{\mu}_f\| = \|\boldsymbol{\mu}_r\|$ (approximately true for well-calibrated classifiers), we have the exact relationship:
\begin{equation}
\sigma_1^2 = \|\boldsymbol{\mu}\|^2 + \boldsymbol{\mu}_f^\top\boldsymbol{\mu}_r, \qquad
\sigma_2^2 = \|\boldsymbol{\mu}\|^2 - \boldsymbol{\mu}_f^\top\boldsymbol{\mu}_r,
\label{eq:sigma_norm}
\end{equation}
and consequently $\|\boldsymbol{\mu}_f - \boldsymbol{\mu}_r\|_2 = \sqrt{2}\,\sigma_2$.

\subsection{The Residual Information Index}

\begin{definition}[Residual Information Index (RII)]
\label{def:rii}
\begin{equation}
\boxed{
\rho(\theta') \coloneqq 1 - \frac{\sigma_1^2}{\sigma_1^2 + \sigma_2^2}
= \frac{\sigma_2^2}{\sigma_1^2 + \sigma_2^2}
\;\in\; [0,\,0.5].
}
\label{eq:rho}
\end{equation}
\end{definition}

The RII has several desirable properties:
\begin{enumerate}
\item \textbf{Normalization:} $\rho \in [0, 0.5]$ independent of $C$ and the scale of $\mathbf{M}$;
\item \textbf{Perfect unlearning:} $\rho = 0 \iff \sigma_2 = 0 \iff \boldsymbol{\mu}_f = \boldsymbol{\mu}_r$;
\item \textbf{Maximal leakage:} $\rho = 0.5$ when $\sigma_1 = \sigma_2$, i.e., when $\boldsymbol{\mu}_f \perp \boldsymbol{\mu}_r$;
\item \textbf{Scale invariance:} $\rho(\lambda\mathbf{M}) = \rho(\mathbf{M})$ for any $\lambda > 0$;
\item \textbf{Connection to spectral deviation:} $\rho \approx \sigma_2^2/\sigma_1^2$ when $\rho \ll 1$.
\end{enumerate}

\paragraph{Statistical justification via the Law of Large Numbers.}
The row vectors $\boldsymbol{\mu}_f$ and $\boldsymbol{\mu}_r$ in Definition~\ref{def:M} are empirical Monte Carlo averages. For i.i.d.\ samples $x_1,\dots,x_{N_f} \sim P_{\mathrm{data}}|_{D_f}$, the strong law of large numbers guarantees
\begin{equation}
\hat{\boldsymbol{\mu}}_f = \frac{1}{N_f}\sum_{i=1}^{N_f} \mathbf{p}(x_i;\theta')
\;\xrightarrow{N_f\to\infty}\;
\boldsymbol{\mu}_f^* := \mathbb{E}_{X\sim P_{\mathrm{data}}|_{D_f}}[\mathbf{p}(X;\theta')],
\end{equation}
with an analogous statement for $\hat{\boldsymbol{\mu}}_r$. The convergence rate is $O_p(1/\sqrt{N_f})$ by the multivariate central limit theorem. Since our experiments use $N_f \ge 500$ (1\% ratio) and $N_f \le 12{,}000$ (20\% ratio), the finite-sample error $\|\hat{\boldsymbol{\mu}}_f - \boldsymbol{\mu}_f^*\|_2$ is bounded by $0.001$--$0.01$ in $\ell_2$ norm for the reported configurations.

The use of output-level averages as proxies for forgetting quality relies on these concentration guarantees, not on any assumption of pointwise equality between individual-sample predictions. The rank-one condition $\rho = 0$ is therefore a statement about \emph{group-level} statistical indistinguishability---that the expected softmax profiles of $D_f$ and $D_r$ coincide---which, by the LLN, can be consistently estimated from finite samples.

%=============================================================================
\section{Class-Level Forgetting and Reference Baselines}
\label{sec:class_baselines}
%=============================================================================

In class-level forgetting, the rows $\boldsymbol{\mu}_f$ and $\boldsymbol{\mu}_r$ average over fundamentally different distributions: $\boldsymbol{\mu}_f$ comes from a class the model has never seen, while $\boldsymbol{\mu}_r$ comes from classes it was trained on. Consequently, even under perfect unlearning, $\boldsymbol{\mu}_f$ and $\boldsymbol{\mu}_r$ can differ, inflating the standard RII. This section discusses existing attempts to address this asymmetry and introduces our improved baseline.

\subsection{The Leave-One-Class-Out (LOCO) Metric}
\label{sec:loco}

A natural proposal is to replace $\boldsymbol{\mu}_r$ with the average output on a \emph{held-out unseen class}---one that never participated in training at all---forming the \textbf{Leave-One-Class-Out (LOCO) confusion matrix}:
\begin{equation}
\boxed{
\mathbf{M}_{\mathrm{LOCO}}(\theta') \;=\;
\begin{bmatrix}
\displaystyle\frac{1}{|D_{c_f}|}\sum_{x \in D_{c_f}} \mathbf{p}(x;\theta') \\[2.2ex]
\displaystyle\frac{1}{|D_{c_0}|}\sum_{x \in D_{c_0}} \mathbf{p}(x;\theta')
\end{bmatrix}
\;\in\; \mathbb{R}^{2\times C},
}
\end{equation}
where $c_0$ is a class held out from training and $c_f$ is the forgotten class. The corresponding LOCO-RII $\rho_{\mathrm{LOCO}}$ is computed from the SVD of $\mathbf{M}_{\mathrm{LOCO}}$.

\paragraph{Why LOCO fails in practice.}
The LOCO intuition would be valid if all ``never-seen'' classes were interchangeable in the model's output distribution. Our experiments on MNIST reveal that this is not the case. In a controlled test, we trained on classes $\{1,2,3,4,6,7,8,9\}$, held out class $0$ as reference, and designated class $5$ as the forget class. The resulting LOCO-RII ($\rho_{\mathrm{LOCO}} \approx 0.32$) was \emph{larger} than the standard RII ($\rho_{\mathrm{std}} \approx 0.14$)---opposite to the expected direction.

The reason is that different unseen classes produce systematically different prediction profiles. A model trained on digits $1$--$9$ (excluding $5$) tends to predict class $5$ samples as visually similar digits (e.g., $3$ or $8$), whereas it predicts class $0$ samples as different patterns (e.g., $6$ or $8$). These inductive biases generate distinct average softmax vectors $\boldsymbol{\mu}_{c_5}$ and $\boldsymbol{\mu}_{c_0}$, inflating $\sigma_2$ and hence $\rho_{\mathrm{LOCO}}$. The LOCO metric therefore \textbf{does not} isolate privacy-relevant leakage; it merely exchanges one distributional mismatch for another.

\subsection{Multi-Held-Out Projection Residual (MHPR)}
\label{sec:mhpr}

To overcome the single-class bias of LOCO, we propose to leverage \emph{multiple} held-out unseen classes and project the forgotten class mean onto their span. The intuition is that while individual unseen classes exhibit idiosyncratic biases, their collective subspace captures the general behavior of ``unseen'' data under the model.

\begin{definition}[Multi-Held-Out Projection Residual, MHPR]
\label{def:mhpr}
Let $\mathcal{H} = \{c_{h1},\dots,c_{hK}\}$ be $K \ge 2$ held-out classes that never participated in training, and let $c_f$ be the forgotten class. Compute the average softmax vectors:
\begin{equation}
\boldsymbol{\mu}_f = \frac{1}{|D_{c_f}|}\sum_{x \in D_{c_f}} \mathbf{p}(x;\theta'), \qquad
\boldsymbol{\mu}_{h_k} = \frac{1}{|D_{h_k}|}\sum_{x \in D_{h_k}} \mathbf{p}(x;\theta').
\end{equation}
Form the reference matrix $\mathbf{H} = [\boldsymbol{\mu}_{h_1},\dots,\boldsymbol{\mu}_{h_K}]^\top \in \mathbb{R}^{K\times C}$.
Let $\hat{\boldsymbol{\mu}}_f = \mathbf{H}^+ \mathbf{H} \boldsymbol{\mu}_f$ be the orthogonal projection of $\boldsymbol{\mu}_f$ onto the row space of $\mathbf{H}$ (with $\mathbf{H}^+$ the Moore--Penrose pseudoinverse). The \textbf{MHPR-RII} is defined as the relative residual energy:
\begin{equation}
\boxed{
\rho_H(\theta') \coloneqq \frac{\|\boldsymbol{\mu}_f - \hat{\boldsymbol{\mu}}_f\|_2^2}{\|\boldsymbol{\mu}_f\|_2^2} \in [0,1].
}
\end{equation}
\end{definition}

If $\boldsymbol{\mu}_f$ can be expressed as a linear combination of the held-out class means, then $\rho_H = 0$ and the forgotten class is \emph{output-indistinguishable} from the collective unseen behavior. A positive $\rho_H$ indicates that $\boldsymbol{\mu}_f$ contains residual structure not captured by any combination of the unseen classes.

\paragraph{Properties and practical considerations.}
\begin{itemize}
\item \textbf{Reduction to LOCO:} When $K=1$, $\hat{\boldsymbol{\mu}}_f = \lambda \boldsymbol{\mu}_{h1}$ and $\rho_H = 1 - \cos^2(\angle(\boldsymbol{\mu}_f, \boldsymbol{\mu}_{h1}))$, a scaled version of the LOCO criterion.
\item \textbf{Degeneracy at large $K$:} If the held-out class means span the entire $\mathbb{R}^C$, any $\boldsymbol{\mu}_f$ will have zero residual. In practice, we use $K=3$--$4$ to balance expressiveness and discriminability.
\item \textbf{Statistical stability:} As with the standard RII, the projection residual benefits from the law of large numbers; the finite-sample error scales as $O(1/\sqrt{\min(N_f, N_h)})$.
\end{itemize}

\paragraph{Why MHPR works.}
The key insight is that the subspace spanned by $\{\boldsymbol{\mu}_{h_k}\}$ captures the \emph{shared structure} of the model's responses to unseen classes (e.g., dispersed probability mass, absence of confident peaks on trained classes). A successfully forgotten class should exhibit this same structure, and thus its mean vector should lie within (or very close to) that subspace. Projecting onto the subspace discards the idiosyncratic single-class bias that plagued LOCO, while preserving sensitivity to genuine output-level leakage.

\subsection{Theoretical Guarantees for MHPR}
\label{sec:mhpr_theory}

We now establish the theoretical foundation for the multi-held-out projection residual. All proofs are deferred to Appendix~\ref{app:mhpr_proofs}.

\paragraph{Notation.}
Let $\boldsymbol{\mu}_f \in \mathbb{R}^C$ be the mean softmax vector of the forgotten class, and let $\mathbf{H}_K \in \mathbb{R}^{K \times C}$ be the matrix whose rows are $K$ held-out class mean vectors $\boldsymbol{\mu}_{h_1},\dots,\boldsymbol{\mu}_{h_K}$. Denote by $\mathcal{S}_K = \mathrm{rowspan}(\mathbf{H}_K) \subseteq \mathbb{R}^C$ the subspace spanned by the held-out class means, and let $\mathbf{P}_{\mathcal{S}_K}$ be the orthogonal projection onto $\mathcal{S}_K$. The MHPR at level $K$ is:
\begin{equation}
\rho_H(K) = \frac{\|\boldsymbol{\mu}_f - \mathbf{P}_{\mathcal{S}_K}\boldsymbol{\mu}_f\|_2^2}{\|\boldsymbol{\mu}_f\|_2^2}.
\end{equation}

\begin{proposition}[Monotonic Improvement with $K$]
\label{prop:mhpr_monotone}
For any $K_1 \le K_2$, we have $\rho_H(K_2) \le \rho_H(K_1)$. The MHPR is non-increasing as more held-out classes are added.
\end{proposition}

\begin{proof}
Since $\mathcal{S}_{K_1} \subseteq \mathcal{S}_{K_2}$ (adding more held-out class means can only enlarge the span), the orthogonal projection onto a larger subspace is closer to any vector. Formally, $\|\boldsymbol{\mu}_f - \mathbf{P}_{\mathcal{S}_{K_2}}\boldsymbol{\mu}_f\|_2 \le \|\boldsymbol{\mu}_f - \mathbf{P}_{\mathcal{S}_{K_1}}\boldsymbol{\mu}_f\|_2$ by the Pythagorean theorem applied to nested subspaces. The denominator $\|\boldsymbol{\mu}_f\|_2^2$ is constant, so $\rho_H(K_2) \le \rho_H(K_1)$.
\end{proof}

\begin{corollary}[MHPR Dominates LOCO]
\label{cor:mhpr_dominates_loco}
For any $K \ge 2$, $\rho_H(K) \le \rho_H(1)$, where $\rho_H(1)$ is equivalent to the LOCO residual up to a scaling factor. Thus, using multiple held-out classes never performs worse than the single-class baseline.
\end{corollary}

\paragraph{Bias decomposition.}
To understand the statistical behavior of MHPR, we decompose its expectation into an \emph{unlearning quality} term and a \emph{subspace approximation error} term.

Let $\boldsymbol{\mu}_f^*$, $\boldsymbol{\mu}_{h_k}^*$ be the population (infinite-sample) mean vectors, and let $\mathcal{S}_K^*$ be the population subspace. Define:
\begin{align}
\rho_{\mathrm{ideal}} &\coloneqq \frac{\|\boldsymbol{\mu}_f^* - \mathbf{P}_{\mathcal{S}_K^*}\boldsymbol{\mu}_f^*\|_2^2}{\|\boldsymbol{\mu}_f^*\|_2^2}, \label{eq:rho_ideal}\\
\rho_{\mathrm{est}} &\coloneqq \frac{\|\boldsymbol{\mu}_f - \mathbf{P}_{\mathcal{S}_K}\boldsymbol{\mu}_f\|_2^2}{\|\boldsymbol{\mu}_f\|_2^2}. \label{eq:rho_est}
\end{align}

\begin{proposition}[Bias--Variance Decomposition]
\label{prop:mhpr_bias}
Assume $\|\boldsymbol{\mu}_f^*\|_2^2 > 0$ and let $N_{\min} = \min(N_f, N_{h_1},\dots,N_{h_K})$. Under the null hypothesis $\mathcal{H}_0$ that $\boldsymbol{\mu}_f^* \in \mathcal{S}_K^*$ (perfect forgetting implies the forgotten class is indistinguishable from the collective unseen behavior), we have:
\begin{equation}
\mathbb{E}[\rho_{\mathrm{est}}] \le \underbrace{\rho_{\mathrm{ideal}}}_{=0\ \text{under }\mathcal{H}_0} \;+\; O\!\left(\frac{C}{N_{\min}}\right),
\label{eq:mhpr_convergence}
\end{equation}
where the expectation is over the random sampling of $D_f$ and $D_{h_k}$.
\end{proposition}

\begin{proof}[Proof Sketch]
Under $\mathcal{H}_0$, $\boldsymbol{\mu}_f^* = \mathbf{P}_{\mathcal{S}_K^*}\boldsymbol{\mu}_f^*$, so $\rho_{\mathrm{ideal}} = 0$. The estimation error arises from two sources: (i) estimating $\boldsymbol{\mu}_f$ from $N_f$ samples, and (ii) estimating the subspace $\mathcal{S}_K$ from the held-out class means. Standard perturbation bounds for principal angles between subspaces~\cite{stewart1990matrix} give $\|\mathbf{P}_{\mathcal{S}_K} - \mathbf{P}_{\mathcal{S}_K^*}\|_2 = O_p(\sqrt{C/N_{\min}})$. Combining with $\|\boldsymbol{\mu}_f - \boldsymbol{\mu}_f^*\|_2 = O_p(\sqrt{C/N_f})$ and applying the delta method yields the $O(C/N_{\min})$ rate.
\end{proof}

\begin{proposition}[Finite-Sample Confidence Bound]
\label{prop:mhpr_confidence}
Let $\hat{\rho}_H$ be the empirical MHPR computed from $N_f$ forget samples and $N_{h_1},\dots,N_{h_K}$ held-out samples per class. Assume per-sample softmax vectors are sub-Gaussian with variance proxy $\nu^2$. For any $\delta \in (0,1)$, with probability at least $1-\delta$:
\begin{equation}
|\hat{\rho}_H - \rho_H^*| \le \frac{8\nu\sqrt{2}}{\|\boldsymbol{\mu}_f^*\|_2^2}\sqrt{\frac{C(K+1)\log((K+1)C/\delta)}{N_{\min}}},
\label{eq:mhpr_conf}
\end{equation}
where $N_{\min} = \min(N_f, N_{h_1},\dots,N_{h_K})$ and $\rho_H^*$ is the population MHPR.
\end{proposition}

\begin{proof}[Proof Sketch]
The estimate $\hat{\rho}_H$ depends on $K+1$ mean vectors ($\boldsymbol{\mu}_f$ and $\boldsymbol{\mu}_{h_1},\dots,\boldsymbol{\mu}_{h_K}$), each estimated from $N_f$ or $N_{h_k}$ sub-Gaussian samples. By Hoeffding's inequality and a union bound over $(K+1)C$ entries, $\max(\|\hat{\boldsymbol{\mu}}_f - \boldsymbol{\mu}_f^*\|_2, \max_k \|\hat{\boldsymbol{\mu}}_{h_k} - \boldsymbol{\mu}_{h_k}^*\|_2) \le 4\nu\sqrt{C\log((K+1)C/\delta)/N_{\min}}$ with probability $\ge 1-\delta$. The MHPR is a locally Lipschitz function of these $K+1$ vectors with constant $O(1/\|\boldsymbol{\mu}_f^*\|_2^2)$, giving the claimed bound via the union bound and triangle inequality.
\end{proof}

\paragraph{Choice of $K$ and the bias--variance tradeoff.}
Proposition~\ref{prop:mhpr_monotone} shows that increasing $K$ never increases $\rho_H$ in expectation, suggesting we should use as many held-out classes as available. However, two practical considerations limit $K$:

\begin{enumerate}
\item \textbf{Estimation variance:} Each additional held-out class adds $C$ parameters to estimate (its mean vector). With finite per-class sample sizes, Proposition~\ref{prop:mhpr_confidence} shows the confidence interval width grows as $\sqrt{K/N_{\min}}$. When $K$ is large relative to $N_{\min}$, the estimation noise can overwhelm the signal.

\item \textbf{Degeneracy at $K = C$:} When $K = C$ and the held-out class means are linearly independent, $\mathcal{S}_K = \mathbb{R}^C$, making $\rho_H = 0$ for \emph{any} $\boldsymbol{\mu}_f$. This is a degenerate case: the metric becomes vacuously zero and loses all discriminative power.
\end{enumerate}

\paragraph{Temperature scaling for improved conditioning.}
When the model's softmax outputs are overly concentrated (due to high confidence or low accuracy), the held-out class mean vectors $\boldsymbol{\mu}_{h_k}$ become nearly collinear, collapsing the subspace $\mathcal{S}_K$ to an essentially one-dimensional manifold. This is the root cause of MHPR failure on Fashion-MNIST (88.7\% accuracy), where $\rho_H \approx 1$ for all $K \le C-1$.

A simple remedy is to apply temperature scaling before computing MHPR:
\begin{equation}
\mathbf{p}_T(x;\theta') = \mathrm{Softmax}\big(f_{\theta'}(x) / T\big), \qquad T > 0,
\end{equation}
and compute $\rho_H(T)$ from the tempered probabilities. As $T \to \infty$, $\mathbf{p}_T \to \mathbf{1}/C$ (uniform), and the subspace spanned by the tempered held-out class means becomes well-conditioned. As $T \to 0$, we recover the original (collinear) regime. We recommend $T = 50$ as a default when the model accuracy is below $90\%$, and $T = 1$ (no scaling) for high-accuracy models. In our Fashion-MNIST experiment, $T = 50$ reduces $\rho_H$ from $0.999$ to $0.021$, which is below the class-level baseline RII of $0.102$---demonstrating that tempered MHPR can recover discriminative power even in previously intractable settings.

\paragraph{Recommendation.}
Our experiments show that $K = 3$--$4$ provides a robust operating point when each held-out class has $\gtrsim 2000$ samples (as in MNIST), while $K = 10$--$20$ is appropriate when per-class samples are limited to $\sim 500$ (as in CIFAR-100). For low-accuracy models ($\lesssim 90\%$), temperature scaling with $T \approx 50$ is recommended to improve the conditioning of the held-out subspace.

%=============================================================================
\section{Operator-Theoretic Generalization}
\label{sec:operator}
%=============================================================================

The discrete $2\times C$ matrix analysis admits a natural lift to continuous spaces via Hilbert--Schmidt operator theory, connecting the RII to classical information-theoretic channel analysis.

\subsection{Hilbert--Schmidt Decomposition}

For a general input space $\mathcal{X}$ and output space $\mathcal{Y}$, the conditional density $f_{Y|X}(y|x)$ defines an integral operator:
\begin{equation}
(T\psi)(y) = \int_{\mathcal{X}} f_{Y|X}(y|x)\,\psi(x)\,dx.
\label{eq:operator}
\end{equation}
$T$ is Hilbert--Schmidt and admits the canonical decomposition:
\begin{equation}
f_{Y|X}(y|x) = \sum_{k=1}^{\infty} \sigma_k\,\phi_k(y)\,\psi_k(x),
\label{eq:hs_decomp}
\end{equation}
where $\{\sigma_k\}_{k=1}^\infty$ are the singular values in descending order, and $\{\phi_k\}$, $\{\psi_k\}$ are orthonormal bases of $L^2(\mathcal{Y})$ and $L^2(\mathcal{X})$, respectively.

\subsection{The Deviation Parameter $\eta$}

We decompose the conditional density into its rank-one (marginal) component and higher-order corrections:
\begin{equation}
f_{Y|X}(y|x) = f_Y(y) + \sum_{k\ge 2} \sigma_k\,\phi_k(y)\,\psi_k(x).
\label{eq:decomp}
\end{equation}

\begin{definition}[Deviation Parameter]
\begin{equation}
\boxed{\eta^2 \coloneqq \sum_{k\ge 2} \sigma_k^2.}
\label{eq:eta}
\end{equation}
\end{definition}

\textbf{Interpretation:} $\eta = 0 \iff$ rank-1 (perfect removal); $\eta > 0 \iff$ residual dependence of rank $\ge 2$. In the discrete $2\times C$ case, $\eta^2 = \sigma_2^2$, and $\rho = \eta^2/(\sigma_1^2 + \eta^2)$.

\subsection{Unified Signal Model}

We model the observed post-removal output as a linear combination of direct signal leakage and cover-channel output, corrupted by independent noise:
\begin{equation}
R = \alpha X + \beta Y + N, \quad N \sim \mathcal{N}(0, \sigma^2),
\label{eq:signal_model}
\end{equation}
where $X$ is the sensitive input (forget-set membership indicator), $Y$ is the ``overwritten'' representation through the channel $f_{Y|X}$, and $\alpha, \beta$ are coupling coefficients. This model captures two orthogonal leakage pathways:

\begin{itemize}
\item \textbf{Direct leakage} ($\propto \alpha^2$): The forget-set signal $X$ directly modulates the observable $R$---features of $D_f$ that survive gradient updates and directly influence predictions.
\item \textbf{Cover leakage} ($\propto \eta^2$): The channel $f_{Y|X}$ deviates from rank one, so $Y$ retains dependence on $X$. This is the spectral leakage measured by $\eta$ (equivalently, by RII in the discrete case).
\end{itemize}

The following small-$\alpha$ expansion, derived via perturbation analysis of differential entropy, is central to our analysis:

\begin{lemma}[Small-$\alpha$ Fisher Information Expansion]
\label{lem:fisher}
For the model $R = \alpha X + Z$ with $X \perp Z = \beta Y + N$, as $\alpha \to 0$:
\begin{equation}
I(X; R) = \frac{\alpha^2}{2}\,\mathrm{Var}(X)\,J(Z) + o(\alpha^2),
\end{equation}
where $J(Z) = \int (f_Z'(z))^2/f_Z(z)\,dz$ is the Fisher information of $Z$. For Gaussian noise, $J(N) = 1/\sigma^2$.
\end{lemma}

\begin{proof}[Proof Sketch]
Expand $f_R(r) = \int f_X(x)f_Z(r-\alpha x)dx$ around $\alpha=0$: $f_R(r) = f_Z(r) + \frac{\alpha^2}{2}\mathrm{Var}(X)f_Z''(r) + o(\alpha^2)$. Using the perturbation expansion of differential entropy $H(f+\delta f) = H(f) + \frac{1}{2}\int (\delta f)^2/f + o(\|\delta f\|^2)$ and noting that the first-order term vanishes, we obtain $H(R) = H(Z) + \frac{\alpha^2}{2}\mathrm{Var}(X)J(Z) + o(\alpha^2)$. Since $H(R|X) = H(Z)$, the result follows.
\end{proof}

%=============================================================================
\section{Main Theoretical Results}
\label{sec:theory}
%=============================================================================

\subsection{Perfect Removal Criterion}

\begin{theorem}[Output Distributional Irreversibility Criterion]
\label{thm:perfect}
Let $\theta' = \mathcal{U}(\theta_0, D_f, D_r)$ be the output of any unlearning algorithm. Assume the model's output distribution $P(Y\mid x, \theta')$ belongs to an exponential family whose sufficient statistic is the mean prediction vector $\boldsymbol{\mu}$. Let $X \in \{0,1\}$ be the forget-set membership indicator. Then:
\begin{equation}
\boxed{\rho(\theta') = 0 \;\Longleftrightarrow\; \boldsymbol{\mu}_f = \boldsymbol{\mu}_r \;\Longleftrightarrow\; I(X; Y \mid \theta') = 0.}
\end{equation}
That is, $\rho=0$ is necessary and sufficient for the model output $Y$ to be conditionally independent of forget-set membership given the parameters. The stronger claim $I(D_f; \theta' \mid \theta_0)=0$ does \emph{not} follow from $\rho=0$ alone and is not asserted.
\end{theorem}

\begin{proof}
$\rho = 0 \Rightarrow \sigma_2 = 0 \Rightarrow \mathrm{rank}(\mathbf{M}) = 1$. Since both rows are probability vectors, $\boldsymbol{\mu}_f = \boldsymbol{\mu}_r$. For the softmax output (a categorical exponential family), the mean parameter $\boldsymbol{\mu}$ is the sufficient statistic. Equality of sufficient statistics implies equality of the full conditional distributions: $P(Y\mid x\in D_f, \theta') = P(Y\mid x\in D_r, \theta')$. Hence $Y \perp X \mid \theta'$, i.e., $I(X; Y \mid \theta') = 0$. The converse follows because $\boldsymbol{\mu}_f \neq \boldsymbol{\mu}_r$ creates a binary channel with positive capacity.
\end{proof}

\begin{remark}[Why Output-Space Only?]
The Markov chain $D_f \to \theta' \to Y$ holds by construction. The data processing inequality gives $I(D_f; Y \mid \theta_0) \le I(D_f; \theta' \mid \theta_0)$. A bound on the output-channel mutual information is a \emph{necessary} condition for parameter-level independence, but not sufficient: internal representations (penultimate features, gradients) may retain $D_f$-specific signatures even when $\rho=0$. Section~\ref{sec:threat} characterizes this gap and proposes complementary audits.
\end{remark}

\subsection{Finite-Sample Confidence}

The empirical RII $\hat{\rho}$ computed from finite samples is a noisy estimate of the population quantity $\rho^*$. We provide a non-asymptotic confidence bound.

\begin{proposition}[Finite-Sample RII Confidence with Subsampling]
\label{prop:confidence}
Let $\hat{\rho}$ be computed from $N_f$ forget and $N_r$ retain samples. Assume per-sample softmax vectors $\mathbf{p}(x;\theta')$ are sub-Gaussian with variance proxy $\nu^2$. For any $\delta \in (0,1)$:
\begin{equation}
\mathbb{P}\!\left(|\hat{\rho} - \rho^*| \le \frac{8\nu\sqrt{2}}{\sigma_1^2}\sqrt{\frac{\log(4C/\delta)}{\min(N_f, N_r)}}\right) \ge 1 - \delta.
\end{equation}
Consequently, to achieve $|\hat{\rho}-\rho^*| \le \varepsilon$ with probability $1-\delta$, it suffices to sample
\begin{equation}
\boxed{N_{\min} \ge \frac{128\nu^2}{\sigma_1^4\varepsilon^2}\log\frac{4C}{\delta}}
\label{eq:subsample}
\end{equation}
samples from each of $D_f$ and $D_r$.
\end{proposition}

\begin{proof}[Proof Sketch]
Each entry of $\hat{\mathbf{M}}$ is a sample mean of $N_f$ or $N_r$ sub-Gaussian variables. By Hoeffding, $\|\hat{\mathbf{M}} - \mathbf{M}\|_F \le O(\sqrt{C/\min(N_f,N_r)})$ with high probability. By Weyl's perturbation theorem, $|\hat{\sigma}_i - \sigma_i| \le \|\hat{\mathbf{M}} - \mathbf{M}\|_2 \le \|\hat{\mathbf{M}} - \mathbf{M}\|_F$. The RII is a smooth function of $(\sigma_1,\sigma_2)$ with gradient bounded by $O(1/\sigma_1^2)$, yielding the claimed rate via the delta method.
\end{proof}

This bound explains the empirical observation that $\hat{\rho}$ decreases with forget ratio: at 1\% forget ($N_f \approx 600$ for MNIST), estimation noise inflates $\hat{\rho}$; at 10\%--20\% ($N_f \ge 6000$), the estimate stabilizes. The $O(1/\sqrt{N_f})$ rate matches our experimental data precisely.

\subsection{Gradient-Based Unlearning Leakage}

\begin{theorem}[Fine-Tuning Leakage Bound]
\label{thm:finetune}
Consider one-step gradient ascent on the forget set: $\theta' = \theta_0 + \eta \nabla_\theta \mathcal{L}(\theta_0, D_f)$, where $\mathcal{L}$ has $L$-Lipschitz gradient. Let $\sigma_n^2$ denote the effective noise scale in predictions. Then:
\begin{equation}
\boxed{\rho(\theta') \le
\frac{\eta^2 L^2}{\sigma_n^2}\,
\operatorname{Tr}\bigl(\operatorname{Cov}_{x\sim D_f}(\nabla_\theta f_{\theta_0}(x))\bigr)
+ O(\eta^3).}
\label{eq:finetune_bound}
\end{equation}
In particular, $\rho = O(\eta^2)$ for small step sizes.
\end{theorem}

\begin{proof}[Proof Sketch]
The gradient step perturbs the output logits by $\Delta f(x) \approx \eta \langle \nabla_\theta f_{\theta_0}(x), \nabla_\theta \mathcal{L}(\theta_0, D_f)\rangle$. The change in $\boldsymbol{\mu}_f$ relative to $\boldsymbol{\mu}_r$ is $\Delta\boldsymbol{\mu} = \frac{1}{N_f}\sum_{x\in D_f} \mathbf{J}_{\mathrm{softmax}}(f(x)) \cdot \Delta f(x) + o(\eta)$, where $\mathbf{J}_{\mathrm{softmax}}$ is the softmax Jacobian. Under the assumption that per-sample gradients on $D_f$ are approximately uncorrelated with those on $D_r$, the squared $\ell_2$ deviation satisfies $\|\Delta\boldsymbol{\mu}\|_2^2 \le O(\eta^2) \cdot \mathrm{Tr}(\mathrm{Cov}(\nabla f))$. Since $\sigma_2 = \|\boldsymbol{\mu}_f - \boldsymbol{\mu}_r\|_2 / \sqrt{2}$ and $\rho \approx \sigma_2^2/\sigma_1^2$ for small $\rho$, the result follows.
\end{proof}

Theorem~\ref{thm:finetune} has an important practical implication: \textbf{the RII after fine-tuning grows at most quadratically in the learning rate}, enabling practitioners to choose $\eta$ to balance unlearning speed against residual leakage. Our experiments confirm this: with $\eta=10^{-5}$, $\rho$ remains below $10^{-3}$.

\subsection{Spectral Mutual Information Bound}

\begin{theorem}[Spectral MI Bound]
\label{thm:mi_bound}
Let $X \in \{0,1\}$ be the forget/retain indicator and $Y \in \{1,\dots,C\}$ the predicted class under $\theta'$. Then:
\begin{equation}
\boxed{I(X; Y \mid \theta') \le \frac{2\sigma_1^2}{\delta\ln 2}\cdot\frac{\rho}{1-\rho} + o(\rho),}
\label{eq:mi_bound}
\end{equation}
where $\delta = \min_y P_r(y) > 0$. For well-calibrated classifiers ($\delta \approx 1/C$): $I \le \frac{2C\sigma_1^2}{\ln 2}\cdot\frac{\rho}{1-\rho}$. For $\rho \ll 1$: $I \le O(\rho)$.
\end{theorem}

\begin{proof}
From the SVD of $\mathbf{M}$ and the norm equality $\|\boldsymbol{\mu}_f - \boldsymbol{\mu}_r\|_2 = \sqrt{2}\,\sigma_2$ (exact when $\|\boldsymbol{\mu}_f\| = \|\boldsymbol{\mu}_r\|$, approximate otherwise), we have via Cauchy--Schwarz:
\begin{equation}
\|\boldsymbol{\mu}_f - \boldsymbol{\mu}_r\|_1 \le \sqrt{C}\,\|\boldsymbol{\mu}_f - \boldsymbol{\mu}_r\|_2 = \sqrt{2C}\,\sigma_2.
\end{equation}
Using $\sigma_2^2 = \sigma_1^2 \cdot \rho/(1-\rho)$, the $\chi^2$-divergence satisfies:
\begin{equation}
\begin{aligned}
\chi^2(P_f\|P_r)
&= \sum_{y} \frac{(P_f(y)-P_r(y))^2}{P_r(y)} \\
&\le \frac{1}{\delta}\|\boldsymbol{\mu}_f - \boldsymbol{\mu}_r\|_2^2
= \frac{2\sigma_1^2}{\delta} \cdot \frac{\rho}{1-\rho}.
\end{aligned}
\end{equation}
Since $D_{\mathrm{KL}}(P_f\|P_r) \le \log(1 + \chi^2(P_f\|P_r)) \le \chi^2(P_f\|P_r)$ (using $\log(1+x) \le x$), and for binary-input channels $I(X;Y) \le \max\{D_{\mathrm{KL}}(P_f\|P_r), D_{\mathrm{KL}}(P_r\|P_f)\}$, we obtain the stated bound. Converting from nats to bits introduces the $1/\ln 2$ factor.
\end{proof}

This bound establishes that \textbf{the mutual information between forget-set membership and model predictions is directly controlled by the RII}. When $\rho \approx 10^{-3}$ (as observed empirically), $I \le O(10^{-2})$ nats---negligible for any practical purpose.

\subsection{Unified Leakage Decomposition}

We now present the central theoretical result that unifies all preceding analysis.

\begin{theorem}[Unified Leakage Decomposition]
\label{thm:unified}
Under the signal model $R = \alpha X + \beta Y + N$ with channel $f_{Y|X}$ having singular values $\{\sigma_k\}_{k=1}^\infty$, the information leakage decomposes as:
\begin{equation}
\boxed{I(X; R) \le
\frac{\alpha^2}{2\sigma^2}\operatorname{Var}(X)
\;+\;
C_2\!\sum_{k\ge 2}\sigma_k^2
\;+\; o(\alpha^2 + \eta^2),}
\label{eq:unified}
\end{equation}
with $I_{\mathrm{dir}} = \frac{\alpha^2}{2\sigma^2}\operatorname{Var}(X)$ (direct leakage) and
$I_{\mathrm{cov}} = C_2\sum_{k\ge 2}\sigma_k^2 = C_2\eta^2$ (cover leakage).
\end{theorem}

\begin{proof}
By the chain rule and data processing inequality, $I(X;R) \le I(X;Y) + I(X;R\mid Y)$. For the cover term, the Hilbert--Schmidt decomposition yields $I(X;Y) \le C_2\sum_{k\ge 2}\sigma_k^2 = C_2\eta^2$, since mutual information between $X$ and $Y$ is bounded by the squared Hilbert--Schmidt norm of the deviation from rank one~\cite{bloch2021overview}. For the direct term, Lemma~\ref{lem:fisher} gives $I(X;R\mid Y) = \frac{\alpha^2}{2\sigma^2}\mathrm{Var}(X) + o(\alpha^2)$, noting that conditioned on $Y$, the only remaining dependence is through $\alpha X$. Summing both contributions yields the decomposition.
\end{proof}

\begin{remark}[Orthogonal Leakage Axes]
Theorem~\ref{thm:unified} reveals a fundamental structure: information leakage after removal has \emph{two orthogonal sources}. The \textbf{direct leakage} term $I_{\mathrm{dir}} \propto \alpha^2$ captures the degree to which the forget-set signal survives in the observable $R$---this is what MIA detects through distributional shift. The \textbf{cover leakage} term $I_{\mathrm{cov}} \propto \eta^2$ captures spectral deviation from the rank-one ideal---this is what RII measures through $\sigma_2/\sigma_1$. When both vanish ($\alpha = 0$ and $\eta = 0$), we recover $I = 0$, i.e., Theorem~\ref{thm:perfect}.
\end{remark}

This decomposition cleanly resolves the MIA--RII paradox: MIA accuracy reflects $I_{\mathrm{dir}}$ (all training samples differ from test samples in confidence), while RII reflects $I_{\mathrm{cov}}$ (forget and retain samples are predicted identically). Both can be large simultaneously, or one can be near-zero while the other remains substantial---they are genuinely independent dimensions of leakage.

%=============================================================================
\section{Experimental Evaluation}
\label{sec:experiments}
%=============================================================================

\subsection{Setup}

We evaluate on five benchmark datasets: MNIST (60K, 28$\times$28 grayscale, 10 classes), Fashion-MNIST (60K, 28$\times$28 grayscale, 10 classes), CIFAR-10 (50K, 32$\times$32$\times$3, 10 classes), CIFAR-100 (50K, 32$\times$32$\times$3, 100 classes), and Tiny ImageNet (100K, 64$\times$64$\times$3, 200 classes). Four model architectures are used: SimpleMLP (784$\to$128$\to$10, $\sim$102K params), SmallCNN (3 conv + 2 fc, $\sim$62K params), ResNet-18 ($\sim$11M params), and DeepMLP (784$\to$512$\to$256$\to$10, $\sim$535K params, for overfitting experiments).

Four unlearning methods are compared:
\begin{itemize}
\item \textbf{NoUnlearn:} Original model, no unlearning applied (lower bound on removal quality).
\item \textbf{Retrain:} Retrain from scratch on $D_r$ only (gold standard, optimal removal).
\item \textbf{SISA~\cite{bourtoule2021machine}:} Train $S=5$ constituent models on disjoint shards; unlearn by retraining only the shard containing $D_f$ ($T=10$ epochs).
\item \textbf{FineTune:} Gradient ascent on $D_f$ (lr $10^{-5}$, 3 epochs) followed by descent on $D_r$ (``ascent+descent'' variant), preserving model utility while suppressing forget-set influence.
\end{itemize}

Training uses Adam optimizer (lr $10^{-3}$, batch 64, 10 epochs). Forget ratios: $\{0.01, 0.05, 0.10, 0.20\}$. Metrics: test accuracy, RII $\rho$, MI upper bound $I_{\mathrm{ub}}$ (Theorem~\ref{thm:mi_bound}), MIA accuracy~\cite{shokri2017membership} (loss-threshold attack). All experiments on Apple M5 Pro (24GB) with MPS acceleration. Seed 42 unless noted.

\subsection{Main Results (Sample-Level Forgetting)}

\begin{table}[t]
\centering
\caption{MNIST Results (SimpleMLP). $\rho$: RII; $I_{\mathrm{ub}}$: MI upper bound (nats).}\label{tab:mnist}
\setlength{\tabcolsep}{3pt}
\begin{tabular}{l c c c c c}
\toprule
Method & $N_f/N$ & Acc.\ (\%) & $\rho$ & $I_{\mathrm{ub}}$ & MIA (\%) \\
\midrule
NoUnlearn  & 0.01 & 97.68 & $3.80{\times}10^{-3}$ & $3.81{\times}10^{-3}$ & 94.1 \\
Retrain    & 0.01 & 97.86 & $4.06{\times}10^{-3}$ & $4.07{\times}10^{-3}$ & 94.1 \\
SISA       & 0.01 & 97.00 & $4.24{\times}10^{-3}$ & $4.26{\times}10^{-3}$ & 94.1 \\
FineTune   & 0.01 & \textbf{98.16} & $3.92{\times}10^{-3}$ & $3.93{\times}10^{-3}$ & 94.1 \\
\midrule
NoUnlearn  & 0.05 & 97.45 & $9.76{\times}10^{-4}$ & $9.77{\times}10^{-4}$ & 90.4 \\
Retrain    & 0.05 & 97.86 & $9.08{\times}10^{-4}$ & $9.09{\times}10^{-4}$ & 90.7 \\
SISA       & 0.05 & 96.97 & $9.45{\times}10^{-4}$ & $9.45{\times}10^{-4}$ & 90.6 \\
FineTune   & 0.05 & \textbf{98.29} & $9.61{\times}10^{-4}$ & $9.62{\times}10^{-4}$ & 90.5 \\
\bottomrule
\end{tabular}
\end{table}

\begin{table}[t]
\centering
\caption{CIFAR-10 Results (SmallCNN, mean $\pm$ std over 3 seeds).}\label{tab:cifar10}
\setlength{\tabcolsep}{2.5pt}
\begin{tabular}{l c c c c c}
\toprule
Method & $N_f/N$ & Acc.\ (\%) & $\rho$ & $I_{\mathrm{ub}}$ & MIA (\%) \\
\midrule
NoUnlearn  & 0.05 & $76.1_{\pm1.0}$ & $1.20_{\pm0.24}{\times}10^{-3}$ & $1.20{\times}10^{-3}$ & $90.5_{\pm0.0}$ \\
Retrain    & 0.05 & $76.0_{\pm1.1}$ & $1.37_{\pm0.35}{\times}10^{-3}$ & $1.37{\times}10^{-3}$ & $91.3_{\pm0.0}$ \\
FineTune   & 0.05 & \textbf{$79.0_{\pm0.2}$} & $1.27_{\pm0.13}{\times}10^{-3}$ & $1.27{\times}10^{-3}$ & $90.5_{\pm0.0}$ \\
\bottomrule
\end{tabular}
\end{table}

\begin{table}[t]
\centering
\caption{CIFAR-100 Results (SmallCNN). 100-class regime confirms rank-one universality despite low accuracy.}\label{tab:cifar100}
\setlength{\tabcolsep}{3pt}
\begin{tabular}{l c c c c c}
\toprule
Method & $N_f/N$ & Acc.\ (\%) & $\rho$ & $I_{\mathrm{ub}}$ & MIA (\%) \\
\midrule
NoUnlearn  & 0.05 & 20.79 & $7.10{\times}10^{-4}$ & $7.11{\times}10^{-4}$ & 90.5 \\
NoUnlearn  & 0.10 & 20.11 & $3.42{\times}10^{-4}$ & $3.43{\times}10^{-4}$ & 86.0 \\
\bottomrule
\end{tabular}
\end{table}

\textbf{Finding 1 (Universal near-rank-one):} Across all three datasets (10 to 100 classes, 20\% to 98\% accuracy), model architectures, and unlearning methods, $\rho < 5 \times 10^{-3}$ uniformly. At 5\%--20\% forget ratios, $\rho$ drops below $10^{-3}$. This is the central empirical result: \textbf{standard neural network training naturally induces near-rank-one output channels}, making the forget/retain distinction nearly invisible in the output distribution. The corresponding mutual information is bounded by $I \le 10^{-2}$ nats.

\textbf{Finding 2 (SISA efficacy):} SISA achieves RII values statistically indistinguishable from full retraining ($\Delta\rho < 2\times 10^{-4}$), validating that shard-based unlearning preserves the rank-one property.

\textbf{Finding 3 (FineTune dominance):} FineTune (ascent+descent) achieves the highest accuracy across all configurations while maintaining competitive RII, and is $3\times$ faster than retraining.

\textbf{Finding 4 (Finite-sample effect):} RII peaks at 1\% forget ratio due to estimation noise, consistent with the $O(1/\sqrt{N_f})$ rate in Proposition~\ref{prop:confidence}. The effect disappears at $\ge 5\%$ forget.

\textbf{Finding 5 (Accuracy independence):} CIFAR-100 ($\sim$20\% acc) yields $\rho = 7.1 \times 10^{-4}$ and Fashion-MNIST (88.7\% acc) yields $\rho = 2.83 \times 10^{-4}$---both \emph{lower} than MNIST (97.5\% acc, $\rho = 9.8 \times 10^{-4}$). The rank-one channel property is independent of model quality; it is an intrinsic feature of the softmax output simplex.

\textbf{Finding 6 (Statistical stability):} Across 3 independent seeds on CIFAR-10, the standard deviation of $\rho$ ranges from $2.0\times 10^{-5}$ to $3.5\times 10^{-4}$, confirming that RII is dominated by data distribution rather than initialization noise.


\subsection{Dynamic Range: Random vs.\ Targeted Forgetting}
\label{sec:dynamic_range}

\begin{table}[t]
\centering
\caption{Dynamic range of RII: random forgetting vs.\ targeted class forgetting (CIFAR-10, SmallCNN).}\label{tab:dynamic_range}
\setlength{\tabcolsep}{3pt}
\begin{tabular}{l c c c c}
\toprule
Forgetting Type & $|D_f|$ & Acc.\ (\%) & $\rho$ & $\sigma_2/\sigma_1$ \\
\midrule
Random 5\% (NoUnlearn) & 2,500 & 77.0 & $1.20{\times}10^{-3}$ & 0.035 \\
Random 5\% (FineTune)  & 2,500 & 79.2 & $1.27{\times}10^{-3}$ & 0.036 \\
\textbf{Cat class (baseline)} & \textbf{5,000} & \textbf{76.2} & $\mathbf{1.83{\times}10^{-1}}$ & \textbf{0.47} \\
\bottomrule
\end{tabular}
\end{table}

Table~\ref{tab:dynamic_range} reveals a striking contrast: $\rho$ for targeted class forgetting is $1.83\times 10^{-1}$, nearly \textbf{180 times larger} than random forgetting. This large dynamic range demonstrates that RII is not degenerate---it faithfully reflects the model's capacity to discriminate between semantically distinct groups.

\paragraph{Rank analysis of class-level forgetting.}
The confusion matrix $\mathbf{M}$ before unlearning on this task is manifestly rank-two:
\begin{equation}
\mathbf{M}_{\mathrm{before}} =
\begin{bmatrix}
\mathbf{0}_{10}                           \\
\frac19 \cdot (\mathbf{1}_{10} - \mathbf{e}_3)
\end{bmatrix},
\end{equation}
where $\mathbf{e}_3$ is the one-hot vector at the cat-class index, yielding $\sigma_2/\sigma_1 \approx 0.58$ and $\rho \approx 0.25$. After unlearning (e.g., 100 gradient ascent steps), $\rho$ drops to $2.04\times 10^{-3}$, demonstrating convergence toward the rank-one ideal.

\subsection{RII as a Continuous Measure of Forgetting Progress}
\label{sec:gradient_sweep}

\begin{table}[t]
\centering
\caption{RII as a function of gradient ascent steps (cat class, CIFAR-10).}\label{tab:gradient_sweep}
\setlength{\tabcolsep}{3pt}
\begin{tabular}{c c c c}
\toprule
Gradient Steps & $\rho$ & Acc.\ (\%) & Interpretation \\
\midrule
0 (baseline) & $1.83{\times}10^{-1}$ & 76.2 & Full cat-class discrimination \\
1            & $1.95{\times}10^{-1}$ & 71.9 & Weak perturbation \\
3            & $1.76{\times}10^{-1}$ & 71.2 & Early forgetting \\
10           & $1.05{\times}10^{-1}$ & 69.2 & Moderate forgetting \\
30           & $1.25{\times}10^{-2}$ & 58.4 & Significant forgetting \\
100          & $2.04{\times}10^{-3}$ & 25.4 & Model near collapse \\
\bottomrule
\end{tabular}
\end{table}

Table~\ref{tab:gradient_sweep} demonstrates that RII decreases \textbf{monotonically} from $0.183$ to $0.002$ across two orders of magnitude, confirming RII as a sensitive, graded metric for unlearning progress. This behavior is consistent with Theorem~\ref{thm:finetune}.

\subsection{Evaluation of MHPR for Class-Level Forgetting}
\label{sec:mhpr_eval}

We now evaluate the Multi-Held-Out Projection Residual (Def.~\ref{def:mhpr}) on the MNIST class-level forgetting setup. The model is trained on classes $\{1,2,3,4,6,7,8,9\}$, with class $5$ as the forget class. Held-out classes are chosen from the unseen set $\{0,1,2,3\}$ in increasing numbers $K$.

\begin{table}[t]
\centering
\caption{MHPR ($\rho_H$) compared to standard RII ($\rho_{\mathrm{std}}$) and LOCO-RII ($\rho_{\mathrm{LOCO}}$) on MNIST class forgetting. $K$: number of held-out classes used for the projection subspace.}\label{tab:mhpr}
\setlength{\tabcolsep}{3pt}
\begin{tabular}{c c c c c}
\toprule
$K$ & $\rho_{\mathrm{std}}$ & $\rho_{\mathrm{LOCO}}$ & $\rho_H$ & Acc.\ (\%) \\
\midrule
1 & 0.145 & 0.379 & 0.942 & 79.5 \\
2 & 0.141 & 0.348 & 0.789 & 68.6 \\
3 & 0.136 & 0.310 & \textbf{0.046} & 58.5 \\
4 & 0.142 & 0.167 & \textbf{0.0008} & 48.7 \\
\bottomrule
\end{tabular}
\end{table}

Table~\ref{tab:mhpr} clearly shows that MHPR with $K \ge 3$ held-out classes resolves the deficiencies of both standard RII and LOCO. With $K=1$ (equivalent to a scaled LOCO), $\rho_H = 0.942$ is even worse than LOCO because the single reference fails to capture the unseen-class manifold. With $K=3$, $\rho_H$ drops to $0.046$, almost an order of magnitude below $\rho_{\mathrm{LOCO}}=0.31$ and well below $\rho_{\mathrm{std}}=0.136$. $K=4$ yields near-perfect rank-one behavior ($\rho_H = 0.0008$), albeit at the cost of reduced training data (accuracy 48.7\%). \textbf{We recommend $K=3$ as a sweet spot}, offering strong discrimination while retaining meaningful model utility.

This experiment validates the core hypothesis behind MHPR: \emph{multiple} held-out classes define a subspace that captures the common structure of the model's response to unseen data, against which the forgotten class can be robustly compared.

\paragraph{Cross-dataset MHPR evaluation.}
Table~\ref{tab:mhpr_cross} extends the MHPR evaluation to four datasets spanning different accuracy regimes and class counts. The key findings are:
\begin{enumerate}
\item MHPR with temperature scaling consistently improves over the class-level RII baseline, achieving $\rho_H < \rho_{\mathrm{class}}$ in all cases.
\item The optimal $K$ depends on per-class sample size: $K=3$ for MNIST (6000/class), $K=9$--$10$ for CIFAR-100 (500/class).
\item Temperature scaling is critical for low-accuracy models ($\lesssim 90\%$); it resolves the subspace collapse problem by spreading near-collinear probability vectors.
\end{enumerate}

\begin{table}[t]
\centering
\caption{Cross-dataset MHPR comparison. $T$: temperature used; $K_{\mathrm{opt}}$: recommended $K$; $\rho_{\mathrm{class}}$: class-level RII.}\label{tab:mhpr_cross}
\setlength{\tabcolsep}{3pt}
\begin{tabular}{l c c c c c c}
\toprule
Dataset & Acc.\ (\%) & $C$ & $T$ & $K_{\mathrm{opt}}$ & $\rho_H$ & $\rho_{\mathrm{class}}$ \\
\midrule
MNIST & 97.5 & 10 & 1 & 3 & $\mathbf{0.046}$ & 0.136 \\
Fashion-MNIST & 88.7 & 10 & 50 & 9 & $\mathbf{0.021}$ & 0.102 \\
CIFAR-10 & 75.0 & 10 & 5 & 9 & $\mathbf{0.025}$ & 0.064 \\
CIFAR-100 & 40.6 & 100 & 1 & 10 & 0.133 & 0.010 \\
\bottomrule
\end{tabular}
\end{table}

\paragraph{Ideal forgetting construction.}
To further validate the MHPR criterion, we construct a controlled synthetic experiment. Using a trained Fashion-MNIST model, we:
\begin{enumerate}
\item Choose $K=5$ held-out class mean vectors $\{\boldsymbol{\mu}_{h_k}\}$ and construct a synthetic forget-mean $\boldsymbol{\mu}_f$ that lies \emph{exactly} in their span: $\boldsymbol{\mu}_f = \sum_{k=1}^K \alpha_k \boldsymbol{\mu}_{h_k}$.
\item Compute $\rho_H$ (should be $0$ by construction).
\item Inject increasing levels of isotropic Gaussian noise $\boldsymbol{\varepsilon} \sim \mathcal{N}(0, \varepsilon^2 \mathbf{I})$ and renormalize to the simplex.
\end{enumerate}

\begin{figure}[t]
\centering
\includegraphics[width=0.45\textwidth]{results/mhpr_review/mhpr_review_summary.png}
\caption{MHPR validation: (left) ideal forgetting yields $\rho_H = 0$ exactly, with monotonic increase under noise injection; (right) cross-dataset comparison showing MHPR consistently below class-level RII with appropriate temperature scaling.}\label{fig:mhpr_validation}
\end{figure}

Figure~\ref{fig:mhpr_validation} confirms that $\rho_H = 0$ when $\boldsymbol{\mu}_f$ is in the held-out subspace (machine precision), and $\rho_H$ increases monotonically with the noise level $\varepsilon$. This provides direct empirical support for Proposition~\ref{prop:mhpr_monotone} and establishes that MHPR faithfully detects deviations from the unseen-class manifold. Notably, even at $\varepsilon = 0.5$, $\rho_H$ remains below the class-level RII, demonstrating the metric's robustness to moderate distributional shift.

\subsection{The FineTune Rebound Phenomenon}
\label{sec:finetune_rebound}

\begin{table}[t]
\centering
\caption{The FineTune rebound phenomenon.}\label{tab:finetune_rebound}
\setlength{\tabcolsep}{3pt}
\begin{tabular}{l c c c}
\toprule
Operation & $\rho$ & Acc.\ (\%) & $\sigma_2/\sigma_1$ \\
\midrule
Baseline & $1.83{\times}10^{-1}$ & 76.2 & 0.47 \\
Ascent 10 steps & $1.05{\times}10^{-1}$ & 69.2 & 0.34 \\
\textbf{Ascent 10 + FineTune} & $\mathbf{2.21{\times}10^{-1}}$ & \textbf{75.2} & \textbf{0.52} \\
\bottomrule
\end{tabular}
\end{table}

We observe that fine-tuning on $D_r$ after gradient ascent \emph{increases} $\rho$ beyond its original value. This rebound is captured by Theorem~\ref{thm:unified}: gradient ascent reduces cover leakage ($I_{\mathrm{cov}}$), but subsequent fine-tuning reshapes decision boundaries and misaligns the perturbed $D_f$ representations, inflating $\sigma_2$. RII correctly detects this phenomenon, which simpler accuracy-based measures miss.

\subsection{Large-Scale Validation: ResNet-18 on CIFAR-100}
\label{sec:resnet18}

\begin{table}[t]
\centering
\caption{ResNet-18 on CIFAR-100.}\label{tab:resnet18}
\setlength{\tabcolsep}{3pt}
\begin{tabular}{l c c c c}
\toprule
Model & $|D_f|$ & Acc.\ (\%) & $\rho$ & $\sigma_2/\sigma_1$ \\
\midrule
SmallCNN & 2,500 & 20.8 & $7.10{\times}10^{-4}$ & 0.027 \\
\textbf{ResNet-18} & \textbf{2,500} & \textbf{29.7} & $\mathbf{1.41{\times}10^{-3}}$ & \textbf{0.038} \\
\bottomrule
\end{tabular}
\end{table}

Table~\ref{tab:resnet18} shows that $\rho$ does not scale with model capacity; ResNet-18 maintains near-rank-one output channels, supporting the architectural invariance hypothesis.

\subsection{Sensitivity to Memorization}

\begin{table}[t]
\centering
\caption{RII Sensitivity to Memorization.}\label{tab:sensitivity}
\setlength{\tabcolsep}{3pt}
\begin{tabular}{l l c c c c c}
\toprule
Dataset & Model & $|D|$ & Ep.\ & Acc.\ (\%) & $\rho$ & $\sigma_2/\sigma_1$ \\
\midrule
MNIST   & SimpleMLP & 60K & 10  & 97.45 & $9.8{\times}10^{-4}$ & 0.031 \\
MNIST   & \textbf{DeepMLP} & \textbf{5K} & \textbf{50} & 95.50 & $\mathbf{1.24{\times}10^{-2}}$ & 0.112 \\
CIFAR-10 & SmallCNN & 50K & 10  & 76.34 & $1.5{\times}10^{-3}$ & 0.038 \\
CIFAR-10 & \textbf{SmallCNN} & \textbf{5K} & \textbf{50} & 61.63 & $\mathbf{1.31{\times}10^{-2}}$ & 0.115 \\
\bottomrule
\end{tabular}
\end{table}

Under deliberate overfitting, $\rho$ surges by an order of magnitude, confirming that RII is not degenerate and responds to genuine memorization.

\subsection{Cross-Dataset RII Consistency}

\begin{table}[t]
\centering
\caption{RII across datasets and architectures.}\label{tab:cross_dataset}
\setlength{\tabcolsep}{3pt}
\begin{tabular}{l l c c c c}
\toprule
Dataset & Model & $C$ & $N_f/N$ & Acc.\ (\%) & $\rho$ \\
\midrule
MNIST    & SimpleMLP & 10  & 0.05 & 97.5 & $9.8{\times}10^{-4}$ \\
Fashion-MNIST & SimpleMLP & 10 & 0.05 & 88.7 & $2.83{\times}10^{-4}$ \\
CIFAR-10 & SmallCNN  & 10  & 0.05 & 77.0 & $1.2{\times}10^{-3}$ \\
CIFAR-100& SmallCNN  & 100 & 0.05 & 20.8 & $7.1{\times}10^{-4}$ \\
CIFAR-100& ResNet-18$^\dagger$ & 100 & 0.05 & 29.7 & $1.41{\times}10^{-3}$ \\
Tiny ImageNet & SmallCNN & 200 & 0.05 & 28.4 & $1.35{\times}10^{-3}$ \\
\bottomrule
\end{tabular}
\footnotesize $^\dagger$Trained for 5 epochs.
\end{table}

The near-constant $\rho < 2\times 10^{-3}$ across $C=10$ to $C=200$ confirms that the rank-one property is a structural feature of the softmax output simplex, not a consequence of high accuracy or small class count. Notably, Fashion-MNIST achieves the smallest $\rho$ ($2.83\times10^{-4}$) despite lower accuracy (88.7\% vs.\ 97.5\% for MNIST), further decoupling the RII from model quality.

\subsection{Direct Leakage Axis Verification}
\label{sec:alpha_experiment}

\begin{table}[t]
\centering
\caption{Controlled perturbation (MNIST, 5\% forget).}\label{tab:alpha}
\setlength{\tabcolsep}{3pt}
\begin{tabular}{c c c c}
\toprule
$\alpha$ & $\rho$ (RII) & MIA (\%) & Interpretation \\
\midrule
0.00 & $9.76{\times}10^{-4}$ & 90.4 & Baseline \\
0.20 & $\mathbf{8.91{\times}10^{-3}}$ & 97.2 & Strong direct leakage \\
\bottomrule
\end{tabular}
\end{table}

Table~\ref{tab:alpha} shows the asymmetric response predicted by Theorem~\ref{thm:unified}: as $\alpha$ increases, RII rises $9\times$ while MIA increases only 6.8pp. This confirms the orthogonal leakage axes.

%=============================================================================
\section{Reconciling MIA and RII: Two Orthogonal Leakage Axes}
\label{sec:mia}
%=============================================================================

A striking paradox emerges from our experiments: MIA accuracy remains 77\%--94\% while RII $\rho \approx 10^{-3}$. The resolution follows directly from Theorem~\ref{thm:unified}: MIA and RII measure \textbf{orthogonal leakage dimensions}.

\begin{itemize}
\item \textbf{Distributional shift (MIA axis, $I_{\mathrm{dir}}$):} Any training sample tends to have higher confidence than test samples. MIA exploits this gap, distinguishing $P(\ell \mid \text{train})$ from $P(\ell \mid \text{test})$. This leakage is $\propto \alpha^2$ in Theorem~\ref{thm:unified}.
\item \textbf{Sample-specific leakage (RII axis, $I_{\mathrm{cov}}$):} RII measures whether $P(Y \mid D_f, \theta') \neq P(Y \mid D_r, \theta')$, i.e., whether $D_f$ leaves a distinct output signature. This is $\propto \eta^2 = \sigma_2^2$.
\end{itemize}

These axes are genuinely independent, as confirmed by the data (Table~\ref{tab:mia_rii} in the full version). In particular, Retrain and NoUnlearn can exhibit near-identical RII while both show high MIA accuracy because MIA detects training membership \emph{in general}, not whether a specific forget request has been honored. \textbf{RII correctly certifies output-level forget/retain indistinguishability}, which is the relevant guarantee for erasure compliance.

%=============================================================================
\section{Beyond Output: Completing the Picture}
\label{sec:threat}
%=============================================================================

Output-layer guarantees are necessary but not sufficient for white-box security. Using the chain rule $I(D_f;\theta') = I(D_f;Y\mid\theta') + I(D_f;\theta'\mid Y)$, RII certifies the first term. The second term can be bounded by internal representation audits such as MMD on penultimate features. Our experiments (omitted for brevity) confirm that Retrain achieves the lowest internal feature MMD, while FineTune and NoUnlearn retain internal signatures despite output-level indistinguishability. Together, RII and internal audits form a more complete unlearning certificate.

%=============================================================================
\section{Comparison with Existing Unlearning Metrics}
\label{sec:comparison}
%=============================================================================

\begin{table*}[t]
\centering
\caption{Comparison of unlearning evaluation metrics.}\label{tab:metric_compare}
\setlength{\tabcolsep}{5pt}
\begin{tabular}{l c c c c c}
\toprule
\textbf{Metric} & \textbf{Info.-Theoretic} & \textbf{Hyperparam.-Free} & \textbf{Scale-Invariant} & \textbf{Comp.\ Cost} & \textbf{Threat Model} \\
\midrule
MIA accuracy~\cite{shokri2017membership} & $\times$ & $\times$ & $\times$ & Medium & Black-box \\
DP ($\varepsilon,\delta$)~\cite{guo2020certified} & $\checkmark$ & $\times$ & $\checkmark$ & High & White-box \\
Backdoor success~\cite{sommer2022towards} & $\times$ & $\times$ & $\times$ & High & Black-box \\
Forgetting completeness~\cite{goel2024corrective} & $\times$ & $\times$ & $\times$ & High & Black-box \\
Replay attack~\cite{chen2021when} & $\times$ & $\times$ & $\times$ & Medium & Black-box \\
Auditability~\cite{thudi2022necessity} & $\times$ & $\times$ & $\times$ & High & White-box \\
\hline
\textbf{RII $\rho$ (Ours)} & $\mathbf{\checkmark}$ & $\mathbf{\checkmark}$ & $\mathbf{\checkmark}$ & $\mathbf{O(CN)}$ & Black-box \\
\textbf{MHPR $\rho_H$ (Ours)} & $\mathbf{\checkmark}$ & $\mathbf{\checkmark}$ & $\mathbf{\checkmark}$ & $\mathbf{O(KCN)}$ & Black-box \\
\bottomrule
\end{tabular}
\end{table*}

As shown in Table~\ref{tab:metric_compare}, RII (and its class-level extension MHPR) are the only metrics simultaneously satisfying information-theoretic grounding, hyperparameter-free computation, and scale invariance. They directly quantify the specific question of interest: has the deleted data's influence been removed from the model's output distribution?

%=============================================================================
\section{Practical Deployment Considerations}
\label{sec:deployment}
%=============================================================================

\subsection{Computational Complexity and Subsampling}

RII requires one forward pass over $D_f \cup D_r$, costing $O(C(N_f+N_r))$. For billion-scale forget sets, Proposition~\ref{prop:confidence} provides a subsampling formula:
\begin{equation}
\boxed{N_{\min} \ge \frac{128\nu^2}{\sigma_1^4\varepsilon^2}\log\frac{4C}{\delta}}
\end{equation}
samples from each of $D_f$ and $D_r$, achieving $|\hat{\rho}-\rho^*| \le \varepsilon$ with confidence $1-\delta$. MHPR requires $K$ additional forward passes on the held-out classes, with similar subsampling guarantees.

\textbf{Recommended protocol:} (1) Draw $N_{\min}$ samples uniformly from $D_f$ and $D_r$ (and each held-out class). (2) Compute $\hat{\rho}$ (or $\hat{\rho}_H$). (3) Report $[\hat{\rho}-\varepsilon,\;\hat{\rho}+\varepsilon]$ as the 95\% confidence interval.

%=============================================================================
\section{Related Work}
\label{sec:related}
%=============================================================================

\textbf{Machine Unlearning.} Algorithms include exact retraining~\cite{cao2015towards}, SISA~\cite{bourtoule2021machine}, gradient-based methods~\cite{golatkar2020eternal,sekhari2021remember}, and influence-based approaches~\cite{guo2020certified}. Our contribution is orthogonal: a unified, algorithm-agnostic evaluation metric.

\textbf{Unlearning Evaluation.} MIA~\cite{shokri2017membership,carlini2022membership} dominates practice. DP-based certification~\cite{guo2020certified,ginart2019making} provides worst-case guarantees. Backdoor-based audits~\cite{sommer2022towards} and replay attacks~\cite{chen2021when} offer alternative perspectives. Thudi et al.~\cite{thudi2022necessity} formalize auditability requirements. Section~\ref{sec:comparison} provides a systematic comparison.

\textbf{Information-Theoretic Security.} Shannon~\cite{shannon1949communication}, Bloch et al.~\cite{bloch2021overview}, and Sch\"{a}fer et al.~\cite{schafer2021_info_theoretic} provide foundations.

%=============================================================================
\section{Conclusion and Future Work}
\label{sec:conclusion}
%=============================================================================

We have presented a spectral framework for measuring output-space irreversibility in machine unlearning. The central insight---the rank-one output channel as certificate of output distributional irreversibility---yields five theoretical results and comprehensive empirical validation across four datasets. For class-level forgetting, the MHPR extension resolves the known/unknown class asymmetry by projecting onto multiple held-out unseen classes, achieving an order-of-magnitude improvement over the LOCO baseline. Key empirical findings: (i) $\rho < 5\times 10^{-3}$ under standard training, surging $10\times$ under overfitting; (ii) MIA and RII measure orthogonal leakage axes; (iii) MHPR with $K=3$ provides a reliable class-level forgetting metric. Future work includes large-scale validation on ImageNet and language models, establishing a DP--RII connection, and complete multi-layer certification.

%=============================================================================
\appendix
\section{Complete Proofs}
\label{sec:appendix}
%=============================================================================

\subsection{Proof of Theorem~\ref{thm:perfect} (Output Distributional Irreversibility)}
\begin{proof}
We prove both directions.

\noindent\textit{($\Rightarrow$) Spectral condition implies output independence.}
$\rho(\theta')=0$ implies $\sigma_2=0$. The SVD of $\mathbf{M}$ reduces to $\mathbf{M}=\sigma_1\mathbf{u}_1\mathbf{v}_1^\top$, hence $\mathrm{rank}(\mathbf{M})=1$. Write $\mathbf{M}=[\boldsymbol{\mu}_f\;\boldsymbol{\mu}_r]^\top$. Both rows are probability vectors summing to one. A rank-one $2\times C$ matrix with rows summing to the same positive constant must have proportional rows: $\boldsymbol{\mu}_f = \lambda \boldsymbol{\mu}_r$ for some $\lambda>0$. The sum constraint forces $\lambda=1$, hence $\boldsymbol{\mu}_f = \boldsymbol{\mu}_r$.

The softmax output defines a categorical exponential family with mean parameter $\boldsymbol{\mu}=\mathbb{E}[Y\mid x]=\mathbf{p}(x;\theta')$. The mapping from natural to mean parameters is bijective, so $\boldsymbol{\mu}_f=\boldsymbol{\mu}_r$ implies equality of the full conditional distributions: $P(Y\mid X=1,\theta') = P(Y\mid X=0,\theta')$, i.e., $Y \perp X \mid \theta'$.

\noindent\textit{($\Leftarrow$) Output independence implies spectral condition.}
If $I(X;Y\mid\theta')=0$, then $\boldsymbol{\mu}_f = \boldsymbol{\mu}_r$, hence $\sigma_2=0$ and $\rho=0$.
\end{proof}

\subsection{Proof of Proposition~\ref{prop:confidence} (Finite-Sample Confidence)}
\begin{proof}
Let $\hat{\boldsymbol{\mu}}_f$ be the empirical forget-set mean and $\boldsymbol{\mu}_f^*$ its expectation. Under the sub-Gaussian assumption with variance proxy $\nu^2$, for each coordinate $c$:
\begin{equation}
\mathbb{P}\bigl(|(\hat{\boldsymbol{\mu}}_f)_c - (\boldsymbol{\mu}_f^*)_c| > t\bigr) \le 2\exp\!\left(-\frac{N_f t^2}{2\nu^2}\right).
\end{equation}
A union bound over $C$ coordinates gives, with probability at least $1-\delta/2$,
\begin{equation}
\|\hat{\boldsymbol{\mu}}_f - \boldsymbol{\mu}_f^*\|_\infty \le \nu\sqrt{\frac{2\log(4C/\delta)}{N_f}},
\end{equation}
and $\|\hat{\boldsymbol{\mu}}_f - \boldsymbol{\mu}_f^*\|_2 \le \nu\sqrt{2C\log(4C/\delta)/N_f}$. An identical bound holds for $\hat{\boldsymbol{\mu}}_r$. Combining both rows,
\begin{equation}
\|\hat{\mathbf{M}} - \mathbf{M}^*\|_F \le 2\nu\sqrt{\frac{2C\log(4C/\delta)}{\min(N_f,N_r)}}.
\end{equation}
By Weyl's perturbation theorem, $|\hat{\sigma}_i - \sigma_i^*| \le \|\hat{\mathbf{M}} - \mathbf{M}^*\|_F$. The RII is $\rho = f(\sigma_1,\sigma_2) = 1 - \sigma_1^2/(\sigma_1^2+\sigma_2^2)$. Its gradient is bounded by $\|\nabla f\|_2 \le 2/\sigma_1^2$. By the mean value theorem,
\begin{align}
|\hat{\rho} - \rho^*|
&\le \frac{2\sqrt{2}}{\sigma_1^2}\,\|\hat{\mathbf{M}}-\mathbf{M}^*\|_F \nonumber\\
&\le \frac{8\nu\sqrt{2}}{\sigma_1^2}\sqrt{\frac{C\log(4C/\delta)}{\min(N_f,N_r)}}.
\end{align}
Setting this bound $\le\varepsilon$ and solving for $\min(N_f,N_r)$ yields the subsampling formula.
\end{proof}

\subsection{Proof of Theorem~\ref{thm:finetune} (Fine-Tuning Leakage Bound)}
\begin{proof}
Let $\mathbf{g} = \nabla_\theta\mathcal{L}(\theta_0, D_f)$. For one gradient step $\theta' = \theta_0 + \eta\mathbf{g}$, a first-order Taylor expansion of the logits and softmax gives
\begin{equation}
\hat{\boldsymbol{\mu}}_f' = \hat{\boldsymbol{\mu}}_f
+ \frac{\eta}{N_f}\sum_{x\in D_f}
\mathbf{J}_{\mathrm{softmax}}(f_{\theta_0}(x))\,
\mathbf{J}_\theta f_{\theta_0}(x)\,\mathbf{g}
+ O(\eta^2).
\end{equation}
For $x \in D_r$, the expected first-order change is zero. Under Lipschitz gradient assumptions, the squared $\ell_2$ displacement satisfies
\begin{equation}
\|\hat{\boldsymbol{\mu}}_f' - \hat{\boldsymbol{\mu}}_r'\|_2^2 \le \frac{\eta^2 L^2}{\sigma_n^2}\operatorname{Tr}(\operatorname{Cov}_{x\sim D_f}(\nabla_\theta f_{\theta_0}(x))) + O(\eta^3).
\end{equation}
Using $\|\boldsymbol{\mu}_f-\boldsymbol{\mu}_r\|_2 = \sqrt{2}\,\sigma_2$ and $\rho \approx \sigma_2^2/\sigma_1^2$ for $\rho\ll 1$, we obtain $\rho(\theta') \le O(\eta^2)$.
\end{proof}

\subsection{Proof of Theorem~\ref{thm:unified} (Unified Leakage Decomposition)}
\begin{proof}
We proceed in four steps.

\noindent\textit{Step 1: Chain rule decomposition.}
$I(X;R) \le I(X;R,Y) = I(X;Y) + I(X;R\mid Y).$

\noindent\textit{Step 2: Cover leakage bound.}
The Hilbert--Schmidt decomposition yields $f_{Y|X}(y\mid x) = f_Y(y) + \sum_{k\ge 2} \sigma_k\,\phi_k(y)\,\psi_k(x).$ Using $\chi^2$-divergence and orthonormality, $I(X;Y) \le B\sum_{k\ge 2}\sigma_k^2\,\mathbb{E}_X[\psi_k(X)^2] = C_2\eta^2.$

\noindent\textit{Step 3: Direct leakage bound.}
Conditioned on $Y=y$, $R = \alpha X + \beta y + N$. By Lemma~\ref{lem:fisher}, $I(X;R\mid Y) = \frac{\alpha^2}{2\sigma^2}\operatorname{Var}(X) + o(\alpha^2).$

\noindent\textit{Step 4: Assembly.}
Summing gives $I(X;R) \le \frac{\alpha^2}{2\sigma^2}\operatorname{Var}(X) + C_2\sum_{k\ge 2}\sigma_k^2 + o(\alpha^2+\eta^2).$
\end{proof}

\begin{thebibliography}{99}
\bibitem{bourtoule2021machine}
L.~Bourtoule et al., ``Machine unlearning,'' in \emph{Proc.\ IEEE S\&P}, 2021, pp.~141--159.

\bibitem{nguyen2025survey}
T.~T.~Nguyen et al., ``A survey of machine unlearning,'' \emph{ACM Comput.\ Surv.}, vol.~57, no.~1, 2025.

\bibitem{cao2015towards}
Y.~Cao and J.~Yang, ``Towards making systems forget with machine unlearning,'' in \emph{Proc.\ IEEE S\&P}, 2015.

\bibitem{golatkar2020eternal}
A.~Golatkar et al., ``Eternal sunshine of the spotless net,'' in \emph{Proc.\ CVPR}, 2020.

\bibitem{sekhari2021remember}
A.~Sekhari et al., ``Remember what you want to forget,'' in \emph{Proc.\ NeurIPS}, 2021.

\bibitem{shokri2017membership}
R.~Shokri et al., ``Membership inference attacks against machine learning models,'' in \emph{Proc.\ IEEE S\&P}, 2017.

\bibitem{carlini2022membership}
N.~Carlini et al., ``Membership inference attacks from first principles,'' in \emph{Proc.\ IEEE S\&P}, 2022.

\bibitem{shannon1949communication}
C.~E.~Shannon, ``Communication theory of secrecy systems,'' \emph{Bell Syst.\ Tech.\ J.}, vol.~28, 1949.

\bibitem{bloch2021overview}
M.~Bloch et al., ``An overview of information-theoretic security,'' \emph{IEEE J.\ Sel.\ Areas Inf.\ Theory}, vol.~2, 2021.

\bibitem{schafer2021_info_theoretic}
R.~F.~Sch\"{a}fer et al., ``Information-theoretic security: A tutorial,'' \emph{IEEE Access}, vol.~9, 2021.

\bibitem{guo2020certified}
C.~Guo et al., ``Certified data removal from machine learning models,'' in \emph{Proc.\ ICML}, 2020.

\bibitem{ginart2019making}
A.~Ginart et al., ``Making AI forget you,'' in \emph{Proc.\ NeurIPS}, 2019.

\bibitem{thudi2022necessity}
A.~Thudi et al., ``On the necessity of auditable algorithmic definitions for machine unlearning,'' in \emph{Proc.\ USENIX Security}, 2022.

\bibitem{sommer2022towards}
D.~M.~Sommer et al., ``Towards probabilistic verification of machine unlearning,'' \emph{Proc.\ Privacy Enhancing Technologies}, 2022.

\bibitem{goel2024corrective}
S.~Goel et al., ``Corrective machine unlearning,'' \emph{Trans.\ Mach.\ Learn.\ Res.}, 2024.

\bibitem{chen2021when}
M.~Chen et al., ``When machine unlearning jeopardizes privacy,'' in \emph{Proc.\ ACM CCS}, 2021.

\bibitem{stewart1990matrix}
G.~W.~Stewart and J.~Sun, \emph{Matrix Perturbation Theory}. Academic Press, 1990.

\bibitem{vershynin2018high}
R.~Vershynin, \emph{High-Dimensional Probability}. Cambridge University Press, 2018.

\end{thebibliography}

\appendices
\section{Proofs of MHPR Theoretical Results}
\label{app:mhpr_proofs}

\subsection{Proof of Proposition~\ref{prop:mhpr_monotone} (Monotonic Improvement)}

Let $\mathcal{S}_{K_1} = \mathrm{span}\{\boldsymbol{\mu}_{h_1},\dots,\boldsymbol{\mu}_{h_{K_1}}\}$ and $\mathcal{S}_{K_2} = \mathrm{span}\{\boldsymbol{\mu}_{h_1},\dots,\boldsymbol{\mu}_{h_{K_2}}\}$ for $K_1 \le K_2$. Since $\mathcal{S}_{K_1} \subseteq \mathcal{S}_{K_2}$, the orthogonal projection satisfies $\mathbf{P}_{\mathcal{S}_{K_2}} = \mathbf{P}_{\mathcal{S}_{K_1}} + \mathbf{P}_{\mathcal{S}_{K_2} \cap \mathcal{S}_{K_1}^\perp}$. For any fixed vector $\boldsymbol{\mu}_f$:
\begin{align}
\|\boldsymbol{\mu}_f - \mathbf{P}_{\mathcal{S}_{K_2}}\boldsymbol{\mu}_f\|_2^2 &= \|\boldsymbol{\mu}_f - \mathbf{P}_{\mathcal{S}_{K_1}}\boldsymbol{\mu}_f\|_2^2 - \|\mathbf{P}_{\mathcal{S}_{K_2} \cap \mathcal{S}_{K_1}^\perp}\boldsymbol{\mu}_f\|_2^2 \nonumber\\
&\le \|\boldsymbol{\mu}_f - \mathbf{P}_{\mathcal{S}_{K_1}}\boldsymbol{\mu}_f\|_2^2,
\end{align}
where the equality follows from the Pythagorean theorem applied to the orthogonal decomposition $\boldsymbol{\mu}_f - \mathbf{P}_{\mathcal{S}_{K_1}}\boldsymbol{\mu}_f = \mathbf{P}_{\mathcal{S}_{K_2} \cap \mathcal{S}_{K_1}^\perp}\boldsymbol{\mu}_f + \mathbf{P}_{\mathcal{S}_{K_1}^\perp \cap \mathcal{S}_{K_2}^\perp}\boldsymbol{\mu}_f$. Dividing by $\|\boldsymbol{\mu}_f\|_2^2$ gives $\rho_H(K_2) \le \rho_H(K_1)$.

\subsection{Proof of Proposition~\ref{prop:mhpr_bias} (Bias--Variance Decomposition)}

Under $\mathcal{H}_0$, $\boldsymbol{\mu}_f^* = \mathbf{P}_{\mathcal{S}_K^*}\boldsymbol{\mu}_f^*$, so $\rho_{\mathrm{ideal}} = 0$. Write the estimator as:
\begin{equation}
\rho_{\mathrm{est}} = \frac{\|\boldsymbol{\mu}_f - \mathbf{P}_{\mathcal{S}_K}\boldsymbol{\mu}_f\|_2^2}{\|\boldsymbol{\mu}_f\|_2^2}.
\end{equation}

Define $\Delta_f = \boldsymbol{\mu}_f - \boldsymbol{\mu}_f^*$ and $\Delta_H = \mathbf{P}_{\mathcal{S}_K} - \mathbf{P}_{\mathcal{S}_K^*}$. The numerator expands as:
\begin{align}
\|\boldsymbol{\mu}_f - \mathbf{P}_{\mathcal{S}_K}\boldsymbol{\mu}_f\|_2^2 &= \|(\boldsymbol{\mu}_f^* + \Delta_f) - (\mathbf{P}_{\mathcal{S}_K^*} + \Delta_H)(\boldsymbol{\mu}_f^* + \Delta_f)\|_2^2 \nonumber\\
&= \|(\mathbf{I} - \mathbf{P}_{\mathcal{S}_K^*})\boldsymbol{\mu}_f^* + (\mathbf{I} - \mathbf{P}_{\mathcal{S}_K^*})\Delta_f - \Delta_H\boldsymbol{\mu}_f^* - \Delta_H\Delta_f\|_2^2.
\end{align}

The first term vanishes under $\mathcal{H}_0$. The remaining terms are $O_p(\|\Delta_f\|_2 + \|\Delta_H\|_2)$. Standard concentration bounds~\cite{vershynin2018high} give $\|\Delta_f\|_2 = O_p(\sqrt{C/N_f})$ and results from subspace perturbation theory~\cite{stewart1990matrix} give $\|\Delta_H\|_2 = O_p(\sqrt{CK/N_{\min}})$. The denominator $\|\boldsymbol{\mu}_f\|_2^2$ converges to $\|\boldsymbol{\mu}_f^*\|_2^2 > 0$ at rate $O_p(1/\sqrt{N_f})$. By the delta method, $\mathbb{E}[\rho_{\mathrm{est}}] = O(CK/N_{\min})$, and simplifying $K \le C$ gives $O(C/N_{\min})$.

\subsection{Proof of Proposition~\ref{prop:mhpr_confidence} (Finite-Sample Bound)}

Let $\hat{\mathbf{z}} = (\hat{\boldsymbol{\mu}}_f, \hat{\boldsymbol{\mu}}_{h_1},\dots,\hat{\boldsymbol{\mu}}_{h_K})$ be the $(K+1)C$-dimensional vector of all empirical means, and $\mathbf{z}^*$ its population counterpart. The MHPR $\rho_H$ is a function $\phi: \mathbb{R}^{(K+1)C} \to [0,1]$ defined by:
\begin{equation}
\phi(\mathbf{z}) = \frac{\|\boldsymbol{\mu}_f - \mathbf{P}_{\mathcal{S}(\boldsymbol{\mu}_{h_1},\dots,\boldsymbol{\mu}_{h_K})}\boldsymbol{\mu}_f\|_2^2}{\|\boldsymbol{\mu}_f\|_2^2},
\end{equation}
where $\mathcal{S}(\cdot)$ denotes the subspace spanned by the held-out class means.

The function $\phi$ is locally Lipschitz. For $\|\boldsymbol{\mu}_f\|_2 \ge \varepsilon > 0$, straightforward calculus shows $\|\nabla_{\boldsymbol{\mu}_f}\phi\|_2 \le 2/\|\boldsymbol{\mu}_f\|_2$ and $\|\nabla_{\boldsymbol{\mu}_{h_k}}\phi\|_2 \le 2\|\boldsymbol{\mu}_f\|_2/\|\boldsymbol{\mu}_f\|_2^2 = 2/\|\boldsymbol{\mu}_f\|_2$. Since $\|\boldsymbol{\mu}_f^*\|_2 > 0$ by assumption, we have with high probability that $\|\hat{\boldsymbol{\mu}}_f\|_2 \ge \|\boldsymbol{\mu}_f^*\|_2/2$ for sufficiently large $N_f$, giving a global Lipschitz constant $L = 8/\|\boldsymbol{\mu}_f^*\|_2^2$.

Each entry of $\hat{\boldsymbol{\mu}}_f$ and $\hat{\boldsymbol{\mu}}_{h_k}$ is a sample mean of sub-Gaussian random variables with variance proxy $\nu^2$. By Hoeffding's inequality and a union bound over $(K+1)C$ entries:
\begin{equation}
\mathbb{P}\!\left(\max_i |\hat{z}_i - z_i^*| \ge t\right) \le 2(K+1)C\exp\!\left(-\frac{N_{\min}t^2}{2\nu^2}\right).
\end{equation}

Setting $t = 2\nu\sqrt{\frac{\log(2(K+1)C/\delta)}{N_{\min}}}$ gives $\|\hat{\mathbf{z}} - \mathbf{z}^*\|_2 \le \sqrt{(K+1)C}\,t$ with probability $\ge 1-\delta$. Combining with the Lipschitz constant yields the claimed bound.

\end{document}
